% Frozen contribution: geometry_frontiers, 29 October 2026.
% Integrate after affine-hko-section.tex. Uses that section's action-body lemma.
% Certificate: exact_triangles_compact.py (Sage QQ/PPL).
\section{A product whose optimal affine position is not Lagrangian}
\label{sec:nonlagrangian-product}

The affine optimality of the pentagon product does not follow from a general
principle reducing products to Lagrangian position. Even a product of two
triangles has an affine image whose systolic ratio exceeds that of every
Lagrangian position of the same product. Here ``Lagrangian position'' means
that both original factor planes become Lagrangian; all independent affine
changes within those planes are allowed.

\begin{theorem}\label{thm:nonlagrangian-triangles}
Let
\[
 \Delta=\{(s,t)\in\mathbb R^2:s\leq1,\ t\leq1,\ -s-t\leq1\},
 \qquad K=\Delta\times\Delta.
\]
Among invertible linear maps $B$ for which both factor planes of $BK$
are Lagrangian, the maximum systolic ratio is $2/3$. The matrix
\[
 A=\begin{pmatrix}
 1&0&0&0\\
 1&-1&1&0\\
 0&1&0&0\\
 -1/4&-1/2&0&-1/4
 \end{pmatrix}
\]
satisfies
\[
 c_{\EHZ}(AK)=\frac{27}{10},\qquad
 \Vol_4(AK)=\frac{81}{16},\qquad
 \rho(AK)=\frac{18}{25}>\frac23.
\]
Consequently, no maximum of the systolic ratio on the full affine orbit
of $K$ occurs in Lagrangian position. By affine equivalence, the same
conclusion holds for the product of any two nondegenerate triangles.
\end{theorem}

The comparison covers \emph{every} Lagrangian position. Both ratios are below
one, so the result gives no Viterbo counterexample. It does not compute the
full affine maximum.

\subsection{Closure and the explicit non-Lagrangian image}

Use the conventions of Lemma~\ref{lem:affine-action-body}. The area of
$\Delta$ is $9/2$, so $\Vol_4(K)=81/4$. Its six support-normalized normals are
\[
 a_0=e_1,\quad a_1=e_2,\quad a_2=-e_1-e_2,\qquad
 a_3=e_3,\quad a_4=e_4,\quad a_5=-e_3-e_4.
\]
Closure forces every feasible weight vector to have the form
\begin{equation}\label{eq:triangle-product-weights}
 \beta_0=\beta_1=\beta_2=\frac t3,\qquad
 \beta_3=\beta_4=\beta_5=\frac{1-t}{3},\qquad 0\leq t\leq1.
\end{equation}
Cyclic shifts of an order do not change its action: the change caused by
moving its first weighted normal to the end is its pairing with twice the
sum of the remaining weighted normals, which vanishes by closure. Thus the
$120$ orders starting with $0$ cover all orders. The endpoints in
\eqref{eq:triangle-product-weights} include smaller supports.

For $x=(1,1,-2,1,1,-1)$, direct multiplication gives
\[
 \Omega=A^{-1}JA^{-\mathsf T},\qquad
 \det A=\frac14,\qquad \mathfrak p(x)=4.
\]
For each order its action is
\[
 Q_\sigma(t)=\frac19\bigl[a t^2+b(1-t)^2+d t(1-t)\bigr],
 \qquad a,b\in\{-1,1\}.
\]
The following finite table gives the greatest mixed coefficient $d$ in
each of the four groups. Each group contains $30$ cyclic orders. A witness
in every row also makes the bound sharp in that group.
\begin{center}
\begin{tabular}{ccclc}
\toprule
$a$&$b$&$\max d$&Attaining order&$\max_{\sigma,t} Q_\sigma(t)$\\
\midrule
$1$&$1$&$4$&$(0,1,3,2,5,4)$&$1/6$\\
$1$&$-1$&$6$&$(0,5,1,3,2,4)$&$5/27$\\
$-1$&$1$&$6$&$(0,3,2,5,1,4)$&$5/27$\\
$-1$&$-1$&$4$&$(0,2,5,1,3,4)$&$1/18$\\
\bottomrule
\end{tabular}
\end{center}
All entries follow by summing the $15$ pairings of the listed six normals;
the accompanying exact verifier enumerates precisely these $120$ orders.
Since $t(1-t)\geq0$, the maximum mixed coefficient bounds the whole group.
The four elementary quadratic maxima occur at $t=1/2,2/3,1/3,1/2$,
respectively. Therefore the full Haim--Kislev maximum is
\[
 h_{\mathcal W_K}(x)=\frac5{27}.
\]
The affine capacity formula gives $c_{\EHZ}(AK)=27/10$, and the volume and
ratio in the theorem follow. The table proves the required global capacity
calculation, rather than merely exhibiting one feasible action.

\subsection{The entire Lagrangian subclass}

For a nondegenerate skew matrix, both diagonal $2\times2$ blocks vanish
if and only if both diagonal blocks of its inverse vanish. Since
$B^{\mathsf T}JB=-\Omega^{-1}$, the two factor planes are Lagrangian exactly
when
\[
 \Omega=\begin{pmatrix}0&M\\-M^{\mathsf T}&0\end{pmatrix},\qquad
 M\in GL(2,\mathbb R).
\]
Conversely, every such $M$ occurs, by taking $B=\operatorname{diag}(M^{-1},I_2)$.
Write $h(M)$ for the action support at this cross form. Homogeneity and
$|\det B|^{-1}=|\det M|$ give
\begin{equation}\label{eq:triangles-cross-ratio}
 \rho(BK)=\frac{|\det M|}{162h(M)^2}.
\end{equation}

The same-factor terms now vanish, so each order has action
$t(1-t)$ times a linear form in $M$. Reversal of an order negates this
linear form. Consequently the global maximum occurs at $t=1/2$, and
$h(M)$ is a finite support function. Set
\[
 C=\{M:h(M)\leq1/12\},\qquad
 u_1=(1,0)^{\mathsf T},\quad u_2=(0,1)^{\mathsf T},\quad
 u_3=(1,-1)^{\mathsf T}.
\]
The second finite identity used in the proof is
\begin{equation}\label{eq:triangles-rankone-hull}
 C=\conv\left\{\pm\frac32u_i u_j^{\mathsf T}:1\leq i,j\leq3\right\}.
\end{equation}
To verify it exactly, form the $120$ action rows at weights $1/6$, multiply
them by $12$, and impose their upper bounds $1$. Rational polyhedral
reduction gives a bounded four-dimensional polytope with $30$ facets and
exactly the $18$ vertices displayed in \eqref{eq:triangles-rankone-hull}.
The verifier asserts equality of these two vertex sets, not an approximate
match. This finite hull identity is the only polyhedral calculation used
in the following determinant bound.

Let $D=\conv\{\pm u_1,\pm u_2,\pm u_3\}$. Grouping any convex combination
in \eqref{eq:triangles-rankone-hull} by its first index writes
\[
 M=\frac32\sum_{i=1}^3u_i v_i^{\mathsf T},\qquad
 v_i\in\alpha_i D,\quad \alpha_i\geq0,\quad\sum_i\alpha_i=1.
\]
Each determinant of two vertices of $D$ has absolute value at most one.
By bilinearity and convexity this implies $|\det(v_i,v_j)|\leq\alpha_i\alpha_j$.
The two-dimensional Cauchy--Binet identity now yields
\[
 |\det M|
 \leq\frac94\sum_{i<j}|\det(u_i,u_j)|\,|\det(v_i,v_j)|
 \leq\frac94\sum_{i<j}\alpha_i\alpha_j
 \leq\frac34.
\]
The last step follows from $\sum_i\alpha_i^2\geq1/3$.
Equality is attained by $\alpha_i=1/3$ and $v_i=u_i/3$, giving
\[
 M_* =\frac12\sum_i u_i u_i^{\mathsf T}
 =\begin{pmatrix}1&-1/2\\-1/2&1\end{pmatrix},\qquad
 \det M_*=\frac34.
\]
It also has $h(M_*)=1/12$: a smaller value would allow dilation within $C$
and contradict the determinant bound. For any nonzero $M$, the matrix
$M/(12h(M))$ belongs to $C$. Equation~\eqref{eq:triangles-cross-ratio} therefore
gives
\[
 \rho(BK)\leq\frac{144}{162}\cdot\frac34=\frac23,
\]
with equality for $M_*$. This proves the sharp global comparison.

Finally, a full affine maximum exists. By
Lemma~\ref{lem:affine-action-body}, the action support is a norm. The continuous
function $|\Pf|$ attains a positive maximum on its compact unit ball. Any
maximizing form is nondegenerate and hence is congruent to the standard
symplectic form; it is realized by an invertible linear image. Since its
ratio is at least $18/25$, no maximizing position has Lagrangian factor planes.

\paragraph{Exact checks and scope.}
The retained verifier
\path{certificates/triangle-position/exact_triangles_compact.py}
uses rational arithmetic
and Sage's PPL backend to reconstruct both finite identities from the six
normals. It verifies the four action-coefficient groups, the explicit matrix,
and the $18$-vertex identity. An independent calculation also checks the
global determinant maximum by enumerating stationary points on faces of
dimensions at most two. The proof above replaces that larger face calculation
by Cauchy--Binet. These checks do not establish priority in the literature.
Published bounds for triangle--convex Lagrangian products concern a different
comparison class. The bounded source comparison is recorded at
\path{reviews/triangle-position-prior-comparison.md}; it does not establish
priority.
